\SourceLeaf{279}%
calculation carried out with ten-digit logarithms. Anyone using this table to compute measurements within this zone may omit some of the last decimal places as they wish if lower precision is sufficient (as it usually is). Whatever form is used to present the \textit{results} of a measurement, consistency requires presenting them with a precision that matches that of the measurements themselves. The figures in the results should always allow the observed quantities to be recovered with the same precision as when they were measured. Thus, if only longitudes and latitudes are chosen for this purpose, even trigonometric measurements of moderate precision would hardly be done justice if the results were given only to the precision with which longitudes and latitudes can be determined astronomically. This would give a false measure of the quality of the work and would often deprive us of its most decisive checks.

\Article{11}
\indent The conformal mapping considered here, from the surface of the ellipsoid to that of the sphere, can be used in more than one way to compute trigonometric measurements. This memoir discusses only its direct use. Other, derived ways of using it for the same purpose are reserved for a second memoir.

The direct use was already sketched in its essentials in the work cited above. A system of triangles formed on the surface of the ellipsoid by shortest lines, or geodesics, is represented on the surface of the sphere by a system of triangles whose angles are exactly equal to the corresponding angles on the spheroid. The sides, unless they are meridian arcs, differ in strict terms from arcs of great circles, but the differences are so small that in most cases the sides may be treated as coinciding with those arcs. At least, where the highest accuracy is required, the deviation can readily be calculated with all the necessary precision, always provided that
